%%%%%%%%%%%%%%%%%%%%%%%%%%%%%%%%%%%%%%%%%%%%%%%%%%%%%%%%%%%%%%%%%%%%%%%%%%%%%%%%%%%%%%%%%%
% Ceci est le template à utiliser pour les milestones des projets d'INFO0947.            %
%                                                                                        %
% Vous devez décommenter et compléter les commandes introduites plus bas (intitule, ...) %
% avant de pouvoir compiler le fichier LaTeX.                                            %
%                                                                                        %
% Le PDF produit doit avoir la forme suivante: grpID_ProjetX_MilestoneY.pdf              %
% où:                                                                                    %
%   - grpID est votre identifiant de groupe (cfr. eCampus)                               %
%   - X est le numéro du projet (1 pour Construction, 2 pour TAD)                        %
%   - Y est le numéro du milestone (1 ou 2)                                              %
%%%%%%%%%%%%%%%%%%%%%%%%%%%%%%%%%%%%%%%%%%%%%%%%%%%%%%%%%%%%%%%%%%%%%%%%%%%%%%%%%%%%%%%%%%

% !TEX root = ./main.tex
% !TEX engine = latexmk -pdf
% !TEX buildOnSave = true
\documentclass[a4paper, 11pt, oneside]{article}

\usepackage[utf8]{inputenc}
\usepackage[T1]{fontenc}
\usepackage[french]{babel}
\usepackage{array}
\usepackage{shortvrb}
\usepackage{listings}
\usepackage[fleqn]{amsmath}
\usepackage{amsfonts}
\usepackage{fullpage}
\usepackage{enumerate}
\usepackage{graphicx}             % import, scale, and rotate graphics
\usepackage{subfigure}            % group figures
\usepackage{alltt}
\usepackage{url}
\usepackage{indentfirst}
\usepackage{eurosym}
\usepackage{listings}
\usepackage{color}
\usepackage[table,xcdraw,dvipsnames]{xcolor}

%%%%%%%%%%%%%%%%% TITRE %%%%%%%%%%%%%%%%
% Complétez et décommentez les définitions de macros suivantes :
\newcommand{\intitule}{Analyse ZigZag : Milestone 1}
\newcommand{\GrNbr}{00} %je ne connaissais pas le numero de mon groupe
\newcommand{\PrenomUN}{Mathias}
\newcommand{\NomUN}{Beckers}
\newcommand{\PrenomDEUX}{Arthur}
\newcommand{\NomDEUX}{Pirick}

%%%%%%%% ZONE PROTÉGÉE : MODIFIEZ UNE DES DIX PROCHAINES %%%%%%%%
%%%%%%%%            LIGNES POUR PERDRE 2 PTS.            %%%%%%%%
\title{INFO0947: \intitule}
\author{Groupe \GrNbr : \PrenomUN~\textsc{\NomUN}, \PrenomDEUX~\textsc{\NomDEUX}}
\date{}
\begin{document}

\maketitle
%%%%%%%%%%%%%%%%%%%% FIN DE LA ZONE PROTÉGÉE %%%%%%%%%%%%%%%%%%%%

%%%%%%%%%%%%%%%% MILESTONE %%%%%%%%%%%%%%%
% Écrivez le contenu de votre milestone ci-dessous.

%%%%%%%%%%%%%%%%%%%%%
\section{Production}

Voici les variables que nous utiliserons :

\begin{itemize}

\item \texttt{int i} : variable d'itération de la boucle, avec $0 < i < N$.

\item \texttt{indice\_start\_valide} : indice de début du plus grand segment valide, 
$0 \le indice\_start\_valide \le N-2$.

\item \texttt{indice\_end\_valide} : indice de fin du plus grand segment valide, 
$1 \le indice\_end\_valide \le N-1$.

\item \texttt{indice\_start\_tempo} : indice de début du segment actuellement vérifié, 
$indice\_end\_valide \le indice\_start\_tempo \le N-2$.

\item \texttt{indice\_end\_tempo} : indice de fin du segment actuellement vérifié, 
$indice\_start\_tempo < indice\_end\_tempo \le N-1$.

\item \texttt{operation\_zigzag} : variable indiquant la relation attendue entre deux
éléments consécutifs du tableau :
\begin{itemize}
\item $0$ : l'élément suivant doit être inférieur au précédent
\item $1$ : l'élément suivant doit être supérieur au précédent
\end{itemize}

\end{itemize}

\subsection*{Spécification de la fonction zigzag}

\begin{verbatim}
/*
@pre  : T initialisé ^ N > 0 ^ i = 0
@post : i = N - 1 ^ T inchangé
*/
\end{verbatim}

Formalisation :

\[
\texttt{analyse\_zigzag(int* T, unsigned int N, unsigned int* start, 
unsigned int* end, unsigned int* length, int* sum)}
\]

\[
\forall i, \ 0 \le i \le N-1 :
\begin{cases}
T[i] > T[i-1] & \text{si } operation\_zigzag = 1 \\
T[i] < T[i-1] & \text{si } operation\_zigzag = 0
\end{cases}
\]

%%%%%%%%%%%%%%%%%%%%%%%%%%%%%%%%%%%%%%%%%%%%%%%%%%%%%%%%%%%%%%%%%%%%%%%%%%%%%%%%%%%%%%%%%%
% Dans cette section, vous rédigez la production demandée pour le milestone.             %
%                                                                                        %
% Soyez le plus précis et complet possible.  Toute production ne correspondant pas à ce  %
% qui est demandé sera considérée comme ``hors sujet'' et la séance de feedback sera     %
% contreproductive.                                                                      %
%%%%%%%%%%%%%%%%%%%%%%%%%%%%%%%%%%%%%%%%%%%%%%%%%%%%%%%%%%%%%%%%%%%%%%%%%%%%%%%%%%%%%%%%%%


\section{Question(s)}
%%%%%%%%%%%%%%%%%%%%%%%
Faut - il écrire dans le programme les predicats intermediaires comme dans la slide 108 du chapitre 2?
%%%%%%%%%%%%%%%%%%%%%%%%%%%%%%%%%%%%%%%%%%%%%%%%%%%%%%%%%%%%%%%%%%%%%%%%%%%%%%%%%%%%%%%%%%
% Dans cette section, vous pouvez rédiger les questions que vous avez sur le projet ou   %
% formuler ce qui reste incompris pour vous.                                             %
%                                                                                        %
% Nous en discuterons lors de la rencontre feedback sur votre production pour le         %
% milestone                                                                              %
%%%%%%%%%%%%%%%%%%%%%%%%%%%%%%%%%%%%%%%%%%%%%%%%%%%%%%%%%%%%%%%%%%%%%%%%%%%%%%%%%%%%%%%%%%

\end{document}
